\chapter{Introduction}
\label{chapter:introduction}

\section{Background and Motivation}

Enterprise infrastructure has grown substantially more complex over the last decade. Kubernetes has become the de-facto standard for container orchestration, but its adoption introduces two interrelated and persistent challenges: \emph{reactive monitoring} and \emph{insufficient network security visibility}.

Traditional threshold-based monitoring systems are designed for static, predictable workloads. In a dynamic Kubernetes environment — where pods are ephemeral, traffic patterns shift continuously, and dozens of microservices interact — fixed alert thresholds produce excessive false positives and frequently miss subtle, multi-dimensional anomalies that precede real outages. The result is alert fatigue for on-call engineers and a fundamentally \emph{reactive} posture: problems are discovered only after they have already impacted services.

On the security side, the default Kubernetes networking model relies on \texttt{iptables} rules managed by \texttt{kube-proxy}. This approach operates exclusively at Layer~3/4 (IP and port), has limited observability, and becomes a performance bottleneck at scale. It offers no visibility into the content or identity of application-layer traffic and provides no mechanism for detecting lateral movement or anomalous flow behaviour within the cluster.

Cilium\footnote{\url{https://cilium.io}} addresses these networking limitations through the use of eBPF (extended Berkeley Packet Filter) --- a kernel-level technology that allows safe, sandboxed programs to run inside the Linux kernel without modifying kernel source code. Cilium enforces network policies at Layer~3 through Layer~7 (including HTTP, gRPC, and Kafka), provides mutual TLS encryption between workloads, and exposes rich, real-time network flow telemetry through its Hubble observability layer.

Critically, the combination of AI-driven monitoring and Cilium-based network security is deeply complementary. Cilium's Hubble component generates granular per-flow telemetry that can serve as a high-value input to machine-learning anomaly detection models. Conversely, AI-detected anomalies in infrastructure metrics can trigger dynamic responses at the network layer. Together these technologies enable an infrastructure that is \emph{self-aware} and \emph{self-defending}: capable of detecting emerging problems proactively and responding to threats autonomously.

\section{Project Scope}

This project investigates how an AI-powered monitoring system and a Cilium-secured Kubernetes network can be designed, integrated, and validated to provide proactive, automated infrastructure management in an enterprise context.

The scope covers:
\begin{itemize}
    \item A study of current Kubernetes monitoring and networking limitations.
    \item The design of an integrated architecture combining Prometheus, an AI anomaly detection engine, Cilium, and Hubble.
    \item A proof-of-concept implementation deployed on a bare-metal Kubernetes cluster.
    \item Empirical evaluation of the solution's effectiveness and performance overhead.
\end{itemize}

\section{Relevance to Enterprise Infrastructure}

Enterprise infrastructure teams are under growing pressure to maintain service reliability while managing increasingly complex, distributed systems. Proactive anomaly detection reduces mean-time-to-detect (MTTD) and mean-time-to-resolve (MTTR) for incidents, directly translating to higher service availability. The shift to zero-trust networking — mandated by frameworks such as NIST SP~800-207 and increasingly required by enterprise security policies — demands the kind of identity-aware, application-layer network policy enforcement that Cilium provides.

By combining these two capabilities, this project addresses a genuine gap in the current enterprise tooling landscape and produces a reusable, open-source-based blueprint that infrastructure teams can adopt.
